\chapter*{总复习：高频辨析与做题模板}
\addcontentsline{toc}{chapter}{总复习：高频辨析与做题模板}

\section*{一、概念定位模板}
\addcontentsline{toc}{section}{一、概念定位模板}
\begin{points}
  \item 问 CPU 给谁用：进程调度。
  \item 问多个进程怎么不抢坏共享变量：同步互斥。
  \item 问互相等资源卡住：死锁。
  \item 问地址怎么变、页表、碎片：存储器管理。
  \item 问缺页、页面置换、抖动：虚拟存储器。
  \item 问文件名、目录、逻辑结构、访问方法：文件系统接口。
  \item 问文件块放磁盘哪里、inode、空闲块：文件系统实现。
  \item 问磁盘臂怎么移动、访问时间：大容量存储。
  \item 问 DMA、缓冲、SPOOLing、设备分配：设备管理。
\end{points}

\section*{二、地址转换题模板}
\addcontentsline{toc}{section}{二、地址转换题模板}
\begin{points}
  \item 看页大小/段大小，确定偏移位数或偏移范围。
  \item 把逻辑地址拆成页号/段号和偏移。
  \item 查页表/段表，判断是否越界或缺页。
  \item 拼物理地址：页式为物理块号 + 页内偏移；段式为段基址 + 段内偏移。
\end{points}

\section*{三、PV 题模板}
\addcontentsline{toc}{section}{三、PV 题模板}
\begin{points}
  \item 找共享资源：需要 mutex 互斥。
  \item 找先后关系：先发生的一方 V，后发生的一方 P。
  \item 找数量约束：空缓冲区用 empty，满缓冲区用 full，资源个数就是信号量初值。
  \item 写顺序：先等条件，再进互斥；先出互斥，再释放条件。
  \item 检查死锁：有没有拿着 mutex 等待别人改变条件。
\end{points}

\section*{四、调度题模板}
\addcontentsline{toc}{section}{四、调度题模板}
\begin{points}
  \item 列表：进程、到达时间、服务时间、优先级/队列。
  \item 找调度点：到达、完成、时间片到、阻塞、抢占、优先级调整。
  \item 画甘特图。
  \item 算完成时间、周转时间、等待时间、带权周转时间。
  \item 抢占式算法每个新进程到达都要判断是否抢占。
\end{points}

\section*{五、页面置换题模板}
\addcontentsline{toc}{section}{五、页面置换题模板}
\begin{points}
  \item 初始化页框为空，按引用串逐个访问。
  \item 命中：不计缺页，但更新 LRU/Clock 等状态。
  \item 缺页：若有空页框直接装入；否则按算法选牺牲页。
  \item 记录每一步页框内容和缺页次数。
  \item 判断 Belady：只针对 FIFO 这类算法可能出现，不要泛化到 LRU/OPT。
\end{points}

\section*{六、文件系统计算题模板}
\addcontentsline{toc}{section}{六、文件系统计算题模板}
\begin{points}
  \item 目录检索题：先算每块可放多少目录项，再算平均需要扫描多少块。
  \item 索引最大文件题：先算每个索引块可放多少块号，再根据索引级数求可寻址数据块数。
  \item 链式分配随机访问题：访问第 $k$ 块通常要从头顺链访问到第 $k$ 块。
  \item 连续分配题：直接定位强，但增长和外部碎片是弱点。
\end{points}

\section*{七、最后一遍背诵关键词}
\addcontentsline{toc}{section}{七、最后一遍背诵关键词}
\begin{multicols}{2}
\begin{points}
  \item OS：资源管理器、控制程序、接口。
  \item 四特征：并发、共享、虚拟、异步。
  \item 接口：命令接口、程序接口、图形接口。
  \item 进程：程序一次执行，PCB 是唯一标志。
  \item 线程：CPU 调度基本单位，共享进程资源。
  \item 调度：高级进系统，中级进内存，低级上 CPU。
  \item 临界区：空闲让进、忙则等待、有限等待、让权等待。
  \item 信号量：P 先减，V 先加。
  \item 死锁：互斥、请求保持、不可剥夺、循环等待。
  \item 分页：逻辑连续、物理离散、页表映射。
  \item 分段：按逻辑结构，便于共享保护。
  \item 虚存：局部性、请求调入、页面置换。
  \item 抖动：缺页太多，系统忙于换页。
  \item 文件：命名信息项序列，目录用于检索。
  \item inode：文件元数据和块指针。
  \item 磁盘：寻道 + 旋转延迟 + 传输。
  \item I/O：DMA、通道、缓冲、SPOOLing。
\end{points}
\end{multicols}
